	\newpage
\section{Ogólne określenie wymagań}		%1
%Określenie celu pracy, co chcemy uzyskać, jakie przewidujemy wyniki













\hspace{0.60cm}Program powinien implementować listę dwukierunkową w języku C++, działającą na stercie, z wykorzystaniem wzorców projektowych Factory i Iterator. Powinien umożliwiać dodawanie, usuwanie, wyświetlanie i czyszczenie elementów listy oraz poruszanie się po niej w obu kierunkach. Każda klasa powinna znajdować się w osobnym pliku, a jej działanie należy przetestować w funkcji main. Projekt powinien zawierać dokumentację wygenerowaną za pomocą Doxygen, dokumentację projektową oraz historię pracy z wykorzystaniem GitHuba.

\subsection{Przykład}  %1.1       

\hspace{0.60cm}Lista dwukierunkowa może być wykorzystana do przechowywania i przeglądania liczb, z możliwością przechodzenia do następnego i poprzedniego elementu. Użytkownik może dodawać i usuwać liczby z wybranego miejsca oraz wyświetlać je w kolejności od początku do końca lub odwrotnie.


\subsection{Instalacja projektu} % 1.2

\vspace{0.3cm} % Lepsza praktyka odstępu niż \hspace

Aby pobrać projekt z serwisu GitHub, należy uruchomić terminal i wpisać następujące polecenie:

\begin{enumerate}
    \item Sklonuj repozytorium:
    \begin{quote}
        \texttt{git clone \href{https://github.com/p-majchrzak/PZ_Projekt_Listy.git}{https://github.com/p-majchrzak/PZ_Projekt_Listy}}
    \end{quote}
\end{enumerate}

 
 
 
 
 